\section{Three Scenarios for the Lending Environment}
\label{sec:scenarios}

% FINAL-RUN DRAFTING NOTES (reviewed 2026-09-21)
% Prose budget: 1,150 words total. Complete only after the final experiment run.
% Table columns: benchmark, bureau visibility, mandatory screening.
% Table rows: beta=0 and beta=1. Each cell gets the same outcome panel: default,
% arrears, adoption, cumulative BNPL volume and quintile breakdowns, with uncertainty.
% State affordability-only scope in the sentence reporting the visibility effect.
% Report the direction and size actually observed; do not assume null policy effects.
% Compare peer and no-peer arms to test whether adoption changes policy responses.
% Cumulative volume distinguishes deferral within the horizon from volume not realised
% within it; it cannot establish permanent abstention beyond the horizon.
% Interpret policy effects alongside the stacking evidence in Section 4. Null effects
% alone do not identify a mechanism. Screening already accounts for visible aggregate
% obligations, even though it imposes no direct facility-count cap.
% Keep cooling-off and facility-cap results in Appendix C, including peer interactions.
% Figure: rq3_interventions_beta1.png with only the retained body scenarios; beta=0
% figure in the appendix, beta=0 rows retained in the body table.
% Float notes: outcome definitions, denominators, replicates, seeds and uncertainty.
% End with an answer to RQ3 supported by the final data; avoid statutory thresholds.

\label{sec:levers}
We address RQ3 through three scenarios. Each regulatory switch is compared with the benchmark, both with and without peer influence:
\begin{enumerate}
\item \textbf{Bureau visibility} ($\mathbb{1}^{\text{bur}}$): \ac{BNPL} obligations reach the bureau and enter the Regulation 23A test, closing the asymmetry of Section~\ref{sec:asymmetry}; the model's representation of the 2026 reporting requirement \cite{businessday_bnpl_2026}, subject to the limit stated there.
\item \textbf{Mandatory affordability screening} ($\mathbb{1}^{\text{aff}}$): the Regulation 23A test is applied to \ac{BNPL} originations themselves, the South African analogue of the \ac{FCA}'s proportionate affordability check \cite{fca_ps26_1, woolard2021}.
\item \textbf{Peer effects} ($\beta > 0$): adoption is socially transmitted through Equation~\ref{eq:peer}. A behavioural mechanism rather than an instrument, it is a scenario because it asks whether the answers to the other two depend on how adoption spreads; its $\beta = 0$ arm is the control of every comparison.
\end{enumerate}
The experimental grid crosses the benchmark, bureau visibility and mandatory screening with $\beta=0$ and $\beta=1$.
This separates each regulatory effect from its interaction with peer adoption.
The reported \ac{NCR} position on incidental credit is relevant to the late fees in Submodel~14 (Section~\ref{sec:background}).
The baseline retains a light origination screen as a modelling assumption; Scenario~2 tests mandatory affordability assessment without assuming that incidental-credit classification alone settles which provisions apply.
We compare cumulative \ac{BNPL} volume as well as default to distinguish borrowing deferred within the horizon from borrowing not undertaken during it.
